\begin{proposition}[深度–高度的双层极值分解：固定深度轨道下高度相位的最强不可辨识与最强可辨识]\label{prop:xi-depth-height-two-level-extremal}
\leanverified{paper\_xi\_depth\_height\_two\_level\_extremal}
固定深度侧 Stieltjes 轨道 \(S(t)=\int\frac{1}{t+a}\,\dd\mu(a)\)（等价固定深度测度 \(\mu\)，从而固定全部衰减幅度/权重口径）。
在该前提下，高度参数 \(\gamma\) 仅能通过格点地址步长 \(\Delta\) 的相位字符 \(e^{-\mathrm{i}\gamma n\Delta}\ (n\in\ZZ)\) 进入可见谱指纹（命题 \ref{con:xi-grid-scan-aliasing-threshold}）。
则存在两条互补的极值原则：
\begin{enumerate}
  \item \textbf{（最强不可辨识）}在单一步长 \(\Delta\) 的制度下，高度参数只能确定至同余群
  \[
  \Gamma_{\Delta}:=\frac{2\pi}{\Delta}\ZZ,
  \]
  且该同余回收不可能通过任何仅依赖深度侧 \(\mu\) 的增强消除（深度侧仅改变量级/衰减而不改变相位不变性群）。
  \item \textbf{（最强可辨识）}若允许引入第二步长 \(\Delta'>0\)，则高度同余群收缩为交
  \(
  \Gamma_{\Delta}\cap\Gamma_{\Delta'}
  \)
  （命题 \ref{con:xi-grid-scan-two-step-dealiasing}），并在 \(\Delta/\Delta'\notin\QQ\) 时达到最强可辨识：\(\Gamma_{\Delta}\cap\Gamma_{\Delta'}=\{0\}\)。
\end{enumerate}
因此，深度轨道与高度载体之间的直积耦合不仅是必要结构（参见深度--高度的最小直积命题），而且具有清晰的极值相图：深度侧任何增强都不能打破高度侧的同余回收，而高度侧的最小增强（双地址）即可把同余回收压缩到最小交。
\end{proposition}

\endinput
